\documentclass[11pt]{article}
\title{Identifiability Limits of a Signal-Memory Hypothesis for the Xenopus laevis Regeneration Switch}
\date{}
\author{SIGMA AI-assisted research synthesis}
\usepackage[T1]{fontenc}
\usepackage[utf8]{inputenc}
\usepackage[margin=1in]{geometry}
\usepackage{amsmath,amssymb,array,booktabs,longtable,pdflscape}
\usepackage{microtype}
\usepackage[hyphens]{url}
\usepackage{titlesec}
\usepackage{tikz}
\usetikzlibrary{fit}
\emergencystretch=2em
\raggedbottom
\linespread{1.07}
\usepackage{xcolor}
\definecolor{blue}{HTML}{24485E}
\titleformat{\section}{\Large\bfseries\sffamily\color{blue}}{\thesection}{0.8em}{}
\titleformat{\subsection}{\large\bfseries\sffamily}{\thesubsection}{0.7em}{}
\setlength{\parskip}{.6em}
\setlength{\parindent}{0pt}
\begin{document}
\maketitle

\begin{abstract}
Does the temporary loss of tail-regeneration competence in Xenopus laevis reflect weakened transmission from early injury signaling to a later regenerative epidermal program? We examine the identifiability of this hypothesis rather than report a discovered mechanism. For unlinked early and late biological units, an animal-level coupling slope is not identifiable: identical marginal distributions admit opposite within-animal associations. This result limits individual prediction without ruling out population-level stage-by-time comparisons. Additional conditional model arguments show why reporter filtering, wound-clock gating, and output turnover complicate interpretations of signal persistence.

Shortened basal-epidermal signal memory remains an untested hypothesis, not a biological refutation or discovery. The available evidence supplies neither verified specimen linkage nor a named intervening state with a validated assay, and it contains no measured biological effect or independent confirmation. The bounded contribution is an analytical framework specifying which observations could distinguish the candidate from restricted alternatives and which missing evidence currently prevents a mechanistic conclusion.
\end{abstract}
\section{Introduction and biological context}
The biological question concerns developmental changes in tail-regeneration competence within Xenopus laevis, including variation between animals. Basal-epidermal regeneration-initiating cells and extracellular-matrix remodeling are treated as prior-art exclusions specified by the project, not as discoveries of this analysis. Their primary sources and subsequent literature have not been inspected, so neither literature novelty nor cross-species generality can be claimed.

The focal hypothesis is that the refractory condition has a shorter-lived basal-epidermal state linking injury input to later regenerative transcription. Under this hypothesis, moving an otherwise matched input nearer the transcriptional response window could produce a larger relative benefit at stage 46 than at stage 40. This proposed timing dependency is narrower than the general proposition that injury signaling or regeneration-initiating cells matter.

The primary estimand is the within-stage relationship between an animal's early wound response and its later epidermal transition, not a difference between developmental-stage averages. Before interpreting that relationship biologically, the analysis asks whether the observations contain the joint information needed to estimate it. Separate stage-by-time distributions can describe population responses without identifying an animal-level temporal coupling rule.

\section{Methods and evidence}
The saved analysis examined conditional counterexamples for unlinked measurements and compared simplified reporter and pulse-response models. Its substantive operations were mathematical constructions and algebraic comparisons, not expression-matrix analysis, biological measurement, or an executed intervention. Directly displayed deductions are distinguished from proof candidates whose complete assumptions or constructions are absent from the bounded extracts.

Evidence eligibility was assessed from the supplied descriptions of GSE245320 and its component series. No original sample manifest, expression matrix, specimen-to-library mapping, or specimen-linked regeneration outcome was available for inspection. Consequently, no replicate counts, pooling structure, longitudinal pairing, assay units, batch adjustment, or animal-level predictive performance were inferred from accession names.

The bounded evidence index omitted 12862 saved research fragments across 17 sources. Claims that could depend on those fragments remain provisional; the full omission ledger is retained with the archived research record.
\section{Question, biological scope, and analytical approach}
The central question is whether a developmental-stage-dependent link between early injury signaling and a later basal-epidermal regenerative program can be identified and distinguished from alternative explanations. The proposed signal-memory hypothesis concerns that temporal dependency, rather than the involvement of an already implicated pathway or cell state. Accordingly, the analysis asks what observations would distinguish weakened injury-to-epidermal transmission from an inability to enter the regenerative state, and what conclusions remain ambiguous under the available measurement structure. Its contribution is conditional analytical discrimination, not an observed mechanism for the Xenopus laevis regeneration switch. [C2, C14]

The biological starting point is the project's proposed comparison of X. laevis tadpoles at developmental stages 40 and 46. The project identifies GSE245320 and its component series\textemdash{}GSE245317 and GSE245318 for bulk RNA sequencing, GSE245312 for single-cell RNA sequencing, and GSE245313 for spatial measurements\textemdash{}as candidate observational resources. The reported bulk collection schedule comprises $0$, $0.5$, and $6$ hours post-amputation and $3$ days post-amputation; these labels describe the proposed comparison, not verified coverage of every stage\textendash{}time combination. Neither accession labels nor the stated schedule establish a complete experimental design. [C1]

The consequential provenance limitation is that original sample records, expression matrices, and specimen-to-library mappings are not present in the supplied evidence. Actual biological units, replicate counts, pooling, early-to-late pairing, tissue boundaries, batches, assay availability, and specimen-linked regeneration outcomes therefore remain unverified; this does not establish that such records are publicly unavailable or that the specimens were unpaired. The intended unit of animal-level inference is the biological replicate, with an independently collected pool, if present, distinguished from an individual animal and technical libraries distinguished from biological replication.

Different accession numbers within the named study are not treated as independent validation. [C1]

The novelty boundary is similarly specific. The project attributes transient basal-epidermal regeneration-initiating cells and extracellular-matrix remodeling defects to the 2024 study identified by DOI 10.1186/s13059-024-03396-3, and treats these topics, together with established signaling candidates, as background rather than proposed discoveries. Because the primary publications and an interaction-specific search of subsequent literature have not been inspected, the timing dependency cannot be described as previously unreported. The biological scope also remains X. laevis: findings in X. tropicalis cannot establish the proposed dependency, and the two-stage comparison cannot by itself explain later recovery of competence. [C2, C1]

The analysis actually supported by the evidence consists of conditional identifiability constructions and simplified dynamical-model comparisons, not an executed transcriptomic reanalysis or biological experiment. Here, identifiability asks whether the specified observations distinguish competing associations or mechanisms, rather than whether a model can be fitted to them. Counterexamples construct different underlying relationships that yield the same permitted observations; dynamical comparisons examine whether a measured temporal response uniquely specifies the biological input or process responsible for it.

Their force depends on explicit assumptions about specimen linkage and measurement dynamics, not on an empirical effect estimated from these accessions. [C14]

This approach separates the evidence needed for a population-level stage-by-time comparison from that needed for an animal-level temporal predictor or a mechanistic interpretation. Linked measurements, unlinked snapshots, and reporter endpoints pose different analytical questions and cannot be substituted for one another merely because they concern the same stages or regenerative program. The resulting argument therefore turns on whether the observations can discriminate the proposed timing dependency from its alternatives, rather than whether known regenerative components differ between developmental conditions. [C1, C2, C14]

\section{Why stage-by-time distributions do not identify animal-level coupling}
Stage-specific early and late response distributions do not, by themselves, identify whether an individual animal\textquoteright{}s injury response predicts its subsequent epidermal transition. This distinction is central to interpreting the proposed signal-memory mechanism: a population-level sequence of responses is not equivalent to evidence that early signaling and later regenerative transcription are coupled within animals. The relevant limitation concerns the joint distribution of those measurements, rather than the precision with which either response is measured separately (C3).

For the proposed Xenopus laevis comparison, let $E_i$ denote animal $i$\textquoteright{}s early wound-response measurement, $R_i$ its later epidermal regeneration-initiating-cell and extracellular-matrix transition measurement, and $S_i$ its developmental stage. The animal is the biological unit for this estimand; paired measurements must therefore refer to the same animal, not merely to samples sharing a stage label. The supplied evidence does not establish specimen linkage, pooling, biological replication, or assay-level sample provenance for GSE245320 and its component series, so the argument below is conditional on early and late observations being unlinked.

Even paired pools would support a pool-level analysis rather than automatically identify animal-level coupling (C3).

A precise target is the within-stage linear-projection slope, $\beta_s=\operatorname{Cov}(E,R\mid S=s)/\operatorname{Var}(E\mid S=s)$, assuming the required moments exist and $\operatorname{Var}(E\mid S=s)>0$. Its numerator asks whether animals with larger early responses also have larger later responses within the same developmental stage. Separate early and late samples can estimate the stage-conditioned marginal distributions of $E$ and $R$, including their means and variances, but do not observe their animal-level covariance. Thus, defining this slope does not make it estimable from unlinked observations, nor does the slope itself constitute a causal measure of signal transmission (C3).

The missing information can be demonstrated exactly. Within a fixed stage, let $U$ be uniformly distributed on $[0,1]$. In Model I, set $E=U$ and $R=U$; in Model II, set $E=U$ and $R=1-U$. Both models give precisely the same uniform marginal distribution for the early response and precisely the same uniform marginal distribution for the later response. Nevertheless, Model I pairs larger early responses with larger later responses, whereas Model II pairs larger early responses with smaller later responses (C3).

In both constructions, $\operatorname{Var}(E)=1/12$, but $\operatorname{Cov}(E,R)=1/12$ in Model I and $\operatorname{Cov}(E,R)=-1/12$ in Model II. Their coupling slopes are consequently $+1$ and $-1$, despite identical observable distributions under independent, unlinked early and late sampling. An estimator receiving only those observations cannot distinguish the models, even if their marginal distributions are known without sampling error.

Assigning positive coupling to stage 40 and negative coupling to stage 46, or reversing those assignments, likewise preserves the stage-specific marginals while reversing the direction of the stage difference in coupling; these are mathematical constructions, not measured expression values or proposed biological trajectories (C3).

Adding more unlinked animals improves estimation of the marginal responses but does not resolve this ambiguity. Additional unlinked sampling times also fail to supply the missing animal correspondence unless justified structural restrictions connect the snapshots. Correlations among stage-by-time averages instead describe relationships among population summaries, not the within-stage animal-level slope defined above. Actual longitudinal links, informative interventions, or explicit identifying restrictions could change that conclusion, but it does not follow from denser cross-sectional sampling alone (C3).

Replacing each animal\textquoteright{}s missing early measurement with its stage\textquoteright{}s early mean creates a separate, exact regression problem. Code stage 40 as $S=0$ and stage 46 as $S=1$, and define the assigned predictor as $X=x_{40}(1-S)+x_{46}S$. Then $SX=x_{46}S$: the proposed interaction is proportional to the stage indicator, while $X$ itself is a linear combination of the intercept and stage. A design matrix containing an intercept, stage, this assigned early-response predictor, and its stage interaction therefore has rank at most two, so an independent coupling interaction cannot be estimated. More biological replicates can improve the precision of the stage means but cannot repair this algebraic redundancy (C4).

These results leave population-level stage-by-time comparisons intact while changing what they can establish about the signal-memory hypothesis. Such comparisons can describe when response distributions differ, but cannot alone show that an animal\textquoteright{}s early response predicts its later epidermal transition or that this association weakens between stages. Evidence for the proposed transmission defect must therefore identify the early\textendash{}late relationship at the intended biological unit, rather than infer it from temporal ordering or stage-assigned averages (C3, C4).

\section{Reporter persistence is not identified signal persistence}
Reporter persistence cannot, by itself, identify the persistence of the biological input proposed to support regeneration. The relevant distinction is between an input acting in an epidermal compartment and a reporter state that accumulates and decays in response to that input. The conditional construction below shows that sustained signaling and an interrupted response with a compensating late pulse can produce identical calibrated readings at $0.5$ and $3$ hours after amputation. This is a measurement-identifiability argument, not an observation of signaling dynamics in Xenopus laevis tadpoles [C5].

Let $t$ denote hours after amputation, $I(t)\geq 0$ the local biological input, and $D(t)$ the reporter state in a fixed epidermal compartment. Within developmental stage $s$, assume $dD/dt=k_sI(t)-\lambda_sD(t)$, with $k_s>0$, $\lambda_s>0$, and $D(0)=0$; the observed signal is $Y(t)=b_s+g_sD(t)$, with known baseline $b_s$ and positive scale $g_s$. These assumptions make the calibrated reporter state exactly observable in the idealized argument, while allowing gain and decay to differ between stages. They exclude uncorrected saturation, changing reporter abundance, transport, bleaching, and additional kinetic states; whether this model describes a usable reporter in wounded tadpole epidermis remains experimentally unvalidated [C5].

Suppressing the stage subscript, the integrating-factor solution is $D(t)=k\int_0^t\exp[-\lambda(t-u)]I(u)\,du$. Consequently, subtracting the predicted carryover of the first reading from the second gives $Q=D(3)-\exp(-2.5\lambda)D(0.5)=k\int_{0.5}^{3}\exp[-\lambda(3-u)]I(u)\,du$. The two readings therefore identify a decay-weighted integral of the intervening input, rather than its temporal arrangement. Later input receives greater weight because less of its reporter contribution has decayed by the terminal observation [C5].

Calibration nevertheless excludes a restricted alternative. Because the input is nonnegative and the integration kernel is strictly positive, $Q=0$ implies that the input is zero almost everywhere between $0.5$ and $3$ hours, whereas $Q>0$ establishes that some input occurred during that interval. A positive carryover-corrected value therefore cannot be explained by residual reporter from the early response alone. It still does not distinguish uninterrupted signaling from renewed input after a gap [C5].

To make that ambiguity explicit, define the sustained candidate by $I_A(t)=1$ throughout $0\leq t\leq3$, using a normalized input scale. Define the interrupted candidate by $I_B(t)=1$ for $0\leq t\leq0.5$, $I_B(t)=0$ for $0.5<t<2.5$, and $I_B(t)=C$ for $2.5\leq t\leq3$, where $C=[1-\exp(-2.5\lambda)]/[1-\exp(-0.5\lambda)]$. Both candidates have the same initial reporter state and identical input histories through the first reading. Their subsequent contributions to $D(3)$ are, respectively, $(k/\lambda)[1-\exp(-2.5\lambda)]$ and $(kC/\lambda)[1-\exp(-0.5\lambda)]$, which are equal by the definition of $C$ [C5].

The terminal equality can also be seen by writing $r=\exp(-\lambda/2)$, so that $C=1+r+r^2+r^3+r^4$ and $1<C<5$. The interrupted candidate then gives $D_B(3)=(k/\lambda)[(1-r)r^5+(1-r^5)]=(k/\lambda)(1-r^6)=D_A(3)$, while both first readings equal $(k/\lambda)(1-r)$. The first term in the interrupted candidate's terminal state is decayed reporter from the early response; the second is the compensating late-pulse contribution. Thus exact knowledge of reporter gain and decay does not make the two endpoint readings sufficient to identify sustained input [C5].

The candidates differ in persistence under an explicit operational definition. If persistence is the longest uninterrupted interval within $[0.5,3]$ on which $I(t)\geq1/2$, the sustained candidate has persistence of $2.5$ hours and the interrupted candidate has persistence of $0.5$ hours. Before the late pulse, their reporter difference is $D_A(t)-D_B(t)=(k/\lambda)[1-\exp(-\lambda(t-0.5))]$ for $0.5<t<2.5$; at $1.5$ hours, this becomes $(k/\lambda)[1-\exp(-\lambda)]>0$. A matched-time intermediate reading therefore discriminates these two specified trajectories, conditional on valid kinetics and sufficient measurement precision [C5].

That separating reading does not identify an unrestricted input history. For a finite sampling schedule, a constant positive input within an unsampled interval can be replaced by a sufficiently amplified pulse near its end that preserves the decay-weighted contribution at the next observation; identical subsequent inputs then preserve all later sampled reporter states. The construction relies on allowing such waveforms without an additional amplitude or rate restriction. In contrast, a continuous noiseless trajectory under the valid model permits recovery of the input almost everywhere through $I(t)=[dD/dt+\lambda D(t)]/k$, although differentiation amplifies measurement noise.

Practical persistence inference therefore requires justified temporal restrictions, uncertainty-aware reconstruction, and a stated minimum detectable gap, rather than additional discrete readings alone [C5].

The counterexample also changes peak input and total exposure, so it is not an exposure-matched intervention isolating a causal effect of persistence. Nor does reporter decay establish the decay rate of an endogenous biological memory state. Its consequence for the regeneration argument is narrower but decisive: a late reporter value conditioned on an early value must be interpreted as filtered input accumulation, not as an identified measure of uninterrupted epidermal signaling. The signal-memory hypothesis therefore requires evidence about the intervening input trajectory beyond the two calibrated endpoints [C5].

\section{What pulse timing can distinguish from signal memory}
Pulse retiming supplies a directional prediction of shortened basal-epidermal signal memory, but a timing-dependent transcriptional response does not by itself identify memory as the defective process. The candidate hypothesis concerns Xenopus laevis tadpoles: an injury-induced state linking basal-epidermal signaling to regenerative transcription would persist for less time at stage 46 than at stage 40. Moving an otherwise matched input pulse closer to the later response window should therefore produce a larger relative transcriptional benefit at stage 46. This is an untested mechanistic prediction, without a measured effect size or established novelty, rather than an observed explanation of the developmental switch [C6].

The simplest terminal-response model makes both the prediction and its ambiguity explicit. Let $s$ denote stage, $p$ pulse time, $A$ effective pulse amplitude, $q_s$ transcriptional inducibility, and $\tau_s$ memory duration; at a terminal wound age of $T=6$ hours, write $\mu_s(p)=q_sA\exp[-(T-p)/\tau_s]$. For two pulse times satisfying $p_2>p_1$, the late-to-early response ratio is $\exp[(p_2-p_1)/\tau_s]$, so $\tau_{46}<\tau_{40}$ predicts a larger relative benefit from retiming at stage 46. This prediction assumes that changing pulse time changes neither effective input amplitude nor inducibility [C6, C7].

At that single terminal time, however, pulse-time-dependent inducibility reproduces the entire response curve exactly. For any positive common memory duration $\tau_0$, define $q_s^*(p)=q_s\exp[-(T-p)(1/\tau_s-1/\tau_0)]$; substitution gives $q_s^*(p)A\exp[-(T-p)/\tau_0]=\mu_s(p)$. Thus, even noiseless observations at many pulse times cannot distinguish stage-dependent memory from common memory accompanied by appropriately varying inducibility. A larger late-pulse benefit would discriminate against a strictly timing-insensitive inducibility defect, but not against inducibility that changes with wound age [C7].

The proposed crossed design addresses this ambiguity by separating pulse onset from elapsed response time. Let $d$ be time since pulse onset and $t=p+d$ absolute wound age; choose two pulse times and two delays such that $p_2-p_1=d_2-d_1=L>0$, with both delays extending beyond a pulse of duration $w$. The resulting schedule contains an overlapping absolute readout time, $p_1+d_2=p_2+d_1$, as well as within-pulse comparisons across delays.

In the proposed terminal-assay implementation, separate randomized animals occupy each design condition, animals are the experimental units, and cells or repeated images remain nested measurements; the contrasts therefore compare condition-level responses rather than paired trajectories within an animal [C8].

These comparisons require a restricted dynamical model, not merely a sampling schedule. Let the incremental intervening state obey $P'=b_su_p-k_sP$, with zero incremental initial state, and let instantaneous nascent transcription be $N=q_sP$ in the memory model. Responses are defined relative to time-matched controls, under linear superposition, positive responses, constant coefficients over the tested window, and no saturation, unmodeled feedback, or hidden parallel memory. The assumptions register in the accompanying table identifies which additional restrictions make the alternative models distinguishable [C8].

Matching input is consequential because equal delivered dose, or equal unweighted intracellular reactive oxygen species exposure, does not guarantee equal effective stimulation of a decaying state. Under a translated finite input profile $u(v)$, the post-pulse state is $P_s(p+d)=C_s\exp(-k_sd)$, where $C_s=b_s\int_0^w\exp(k_sv)u(v)\,dv$. The exponential weighting means that redistributing the same exposure within the pulse can change the encoded amplitude. The timing comparison therefore requires matched time-resolved, compartment-specific effective input profiles and the stipulated encoding behavior, rather than area matching alone [C6, C8].

Within those restrictions, comparing delays within each pulse arm removes pulse-specific amplitude. Writing $N_{sij}=N_s(p_i+d_j)$, the memory model gives $\log(N_{si2}/N_{si1})=-k_sL$, whereas a pulse-time-only encoding model with common post-pulse decay gives $\log(N_{si2}/N_{si1})=-k_0L$. Pulse-time-dependent encoding can therefore explain a terminal timing benefit without predicting unequal stage-specific delay slopes in this restricted comparison. Equal slopes, constant encoding gates, or inadequate precision would be nondiscriminating outcomes rather than positive evidence for either mechanism [C8].

The shared absolute readout time supplies a different cancellation, directed at a pure wound-clock gate. Under the restricted output-gate model, $N_{sij}=q_sC_sg_s(p_i+d_j)\exp(-k_0d_j)$, where encoding is independent of pulse time and the gate depends only on absolute wound age. At the overlap, the gate cancels and $K_s=\log(N_{s21}/N_{s12})/L=k_0$, whereas the memory model gives $K_s=k_s$; shortened stage-46 memory therefore predicts $K_{46}-K_{40}>0$, while the common-memory output-gate model predicts zero. Without an overlapping wound age, an unrestricted positive gate could fit the separate readout times, so crossing pulse times and delays alone would not provide this discrimination [C8].

The strongest remaining alternative combines pulse-time encoding with readout-time gating, and it defeats the output-only interpretation exactly. Consider $N=a_s(p)g_s(t)\exp(-k_0d)$, with $a_s(p)=q_sC_s\exp[(k_s-k_0)p]$ and $g_s(t)=\exp[-(k_s-k_0)t]$. Because $t=p+d$, substitution yields $N=q_sC_s\exp[(k_s-k_0)p-(k_s-k_0)(p+d)-k_0d]=q_sC_s\exp(-k_sd)$, identical to the memory-model response. This equivalence holds across pulse times and delays, not just at the four proposed observations, and shows why the restricted cancellations cannot establish memory against simultaneous encoding and wound-clock effects [C9].

Accumulated transcription introduces a separate ambiguity even when encoding and readout gates are held constant. The reported conditional model argument couples the intervening state to nascent transcription and an accumulated output satisfying $R'=N-\lambda_sR$; under zero-state linear dynamics, exchanging signal-state decay and output turnover leaves the impulse response unchanged. This exchange-symmetry argument should be retained as a conditional identification limit, not a general theorem about biological transcription.

It identifies the separating information as a calibrated trajectory of the actual intervening state together with nascent transcription, or independently calibrated output turnover, rather than accumulated output alone [C10].

Pulse timing consequently tests a temporal-transfer explanation more sharply than a single endpoint does, but only after the competing models have been restricted in biologically defensible ways. The crossed comparisons can challenge pulse-time-only encoding and readout-time-only gating separately; their success cannot exclude the exact combined-gate construction or assign accumulated-output persistence to signal memory. The central argument therefore changes from treating timing-dependent rescue as evidence of shortened memory to treating it as a conditional signature whose mechanistic interpretation requires information beyond transcriptional output [C7\textendash{}C10].

\begin{quote}
\textbf{Diagram (readable relationship list).} The candidate signal-memory mechanism and its observation-level alternatives are hypotheses and mathematical models, not measured biological relationships. A persistent reporter, accumulated output, or pulse-timing response cannot by itself identify persistence of the unnamed intervening state.
\begin{description}
\item[\textbf{Node 1.}] Basal-epidermal injury input
\item[\textbf{Node 2.}] Possible pulse-time-dependent encoding
\item[\textbf{Node 3.}] Unnamed intervening state $P$; candidate stage-dependent decay
\item[\textbf{Node 4.}] Regenerative transcriptional response
\item[\textbf{Node 5.}] Possible absolute-wound-age gate
\item[\textbf{Node 6.}] Accumulated output with turnover
\item[\textbf{Node 7.}] Reporter accumulation and decay
\item[\textbf{Node 8.}] Observed endpoint or trajectory
\item[\textbf{Edge 1.}] Node 1 points to node 2: Alternative model permits time-dependent encoding
\item[\textbf{Edge 2.}] Node 2 points to node 3: Hypothesized state induction
\item[\textbf{Edge 3.}] Node 3 points to node 4: Hypothesized coupling
\item[\textbf{Edge 4.}] Node 5 points to node 4: Alternative model permits wound-age-dependent gating
\item[\textbf{Edge 5.}] Node 4 points to node 6: Accumulation can mix induction and turnover
\item[\textbf{Edge 6.}] Node 1 points to node 7: Reporter-model input
\item[\textbf{Edge 7.}] Node 7 points to node 8: Filtered signal measurement
\item[\textbf{Edge 8.}] Node 6 points to node 8: Accumulated-output measurement
\item[\textbf{Boundary 1.}] Unvalidated biological hypothesis and competing gates: node 1, node 2, node 3, node 4, node 5
\item[\textbf{Boundary 2.}] Measurement and output-dynamics alternatives: node 6, node 7, node 8
\end{description}
\par\smallskip{\footnotesize Evidence: [Research session 22, artifact 3], [Research session 22, artifact 2], [Research session 22, artifact 5], [Research session 22, artifact 6]; claims: C5, C6, C7, C9, C10, C11}
\end{quote}
\begin{longtable}{@{}>{\raggedright\arraybackslash}p{\dimexpr(\linewidth-2\tabcolsep*2)/111*27\relax}>{\raggedright\arraybackslash}p{\dimexpr(\linewidth-2\tabcolsep*2)/111*42\relax}>{\raggedright\arraybackslash}p{\dimexpr(\linewidth-2\tabcolsep*2)/111*42\relax}@{}}
\caption{Assumptions register for interpreting pulse-time and delay comparisons. These are proposed model restrictions, not measured properties of Xenopus laevis basal epidermis.}\\
\toprule
\textbf{Restricted model} & \textbf{Location of the proposed stage dependence} & \textbf{Restriction required for the timing interpretation}\\
\midrule
\endfirsthead
\toprule
\textbf{Restricted model} & \textbf{Location of the proposed stage dependence} & \textbf{Restriction required for the timing interpretation}\\
\midrule
\endhead
\midrule
\multicolumn{3}{r}{\footnotesize Continued on next page}\\
\endfoot
\bottomrule
\endlastfoot
Signal memory & Decay rate $k_s$ of the intervening state & Nascent transcription satisfies $N=q_sP$, with $q_s$ independent of pulse time and absolute wound age; translated inputs yield the same encoded amplitude across pulse times within stage.\\
Pulse-time encoding only & Encoded amplitude $C_s(p)$ & Post-pulse decay is the common rate $k_0$; the encoding gate does not create continuing unmeasured input after pulse termination.\\
Readout-clock gating only & Transcriptional gate $g_s(t)$ & Post-pulse decay is the common rate $k_0$; encoded amplitude is pulse-time independent, and the gate depends only on absolute wound age, not pulse history.\\
Accumulated-output turnover & Turnover rate $\lambda_s$ of accumulated output & Signal decay and transcriptional gain are otherwise common or constant as specified; interpreting output persistence additionally requires separating turnover from intervening-state decay.\\
\end{longtable}
\par\smallskip{\footnotesize Evidence: [Research session 22, artifact 5], [Research session 22, artifact 6]; claims: C8, C10}
\section{Operational requirements and bounded conclusions}
The decisive requirement for testing shortened signal memory is an identifiable endogenous state between injury input and regenerative output, not another measurement of either endpoint. For the proposed Xenopus laevis stage-40 versus stage-46 comparison, no molecular basal-epidermal state corresponding to the intervening variable $P$ has been identified in the supplied evidence. Defining $P$ by its dynamical role therefore does not yet provide something that can be measured or selectively perturbed.

An upstream reactive oxygen species sensor would measure input unless a distinct intervening role were demonstrated, while a downstream transcriptional reporter would retain the measurement-persistence ambiguity; neither can simply be renamed as the missing state [C11].

A contingent validation sequence would begin with a named endogenous candidate and evidence establishing its localization in the relevant basal epidermis. Candidate-specific measurements would then require specificity, adequate temporal resolution, stable normalization, and evidence that measurement does not materially perturb the response. Measurement kinetics must be calibrated independently within each developmental stage and relevant injury context, rather than estimated freely from the same trajectories used to infer biological decay. An orthogonal assay of the endogenous state would be required to corroborate the reconstructed trajectory before interpreting a stage difference as altered persistence [C13].

Verified cessation of effective input is equally necessary. Ending stimulus delivery or washing the external medium would not, without intracellular measurements or independently justified bounds, establish that the candidate is no longer being driven. Exposure validation would need to address the achieved local profile, including peak, duration, shape, localization, and interactions with endogenous signaling; equal delivered dose or equal integrated exposure alone would not establish equivalent effective input. These requirements distinguish a decay experiment from an experiment in which continued stimulation produces an apparently persistent state [C13].

Future timing comparisons must also distinguish pulse timing after amputation, elapsed time since the pulse, and absolute wound age. Measurements at matched post-pulse delays are needed to compare response kinetics, while measurements at matched wound ages are needed to assess dependence on the developmental course of the wound. Neither comparison alone substitutes for the other, and both require validated exposure profiles and contemporaneous controls. Noninjurious delivery, sham and vehicle controls, and assessment of wound size, closure, viability, and tissue composition would help constrain explanations based on treatment injury or altered wound progression rather than the designated timing dependency [C13].

Only after these conditions are satisfied would corrected candidate-state decay provide a discriminating observation. Faster decay at stage 46 than at stage 40 after matched effective input has ended, corroborated by an orthogonal endogenous-state assay, would support a candidate-specific persistence difference over the restricted explanation of common biological decay with stage-dependent measurement kinetics. Conversely, a sufficiently precise comparison excluding the prespecified faster-decay effect, including equivalence within a justified interval, would oppose that designated hypothesis.

A nonsignificant but imprecise difference would remain inconclusive, and neither outcome would eliminate unrestricted hidden-state or combined-gating explanations [C13].

A persistence difference would still be distinct from causal mediation of regeneration. A subsequent intervention would need to change the candidate\textquoteright{}s persistence while separating that change from achieved exposure, toxicity, composition, pulse-time inducibility, and direct effects on the later epidermal program. Testing a downstream bypass alongside an early-signal perturbation would help distinguish impaired transmission from an inability to activate the downstream regenerative state.

Even an apparent pulse rescue would not by itself identify the intervening mechanism, and evidence against a practically useful rescue would require uncertainty bounds that exclude a prespecified useful benefit after the intended persistence change has been verified [C13].

The biological unit and sample provenance must be established before these observations can support individual-level inference. For GSE245320 and its named component series, the supplied evidence does not verify specimen-to-library mappings, animals versus pools, temporal pairing, replication, batch structure, or specimen-linked regeneration outcomes. These relationships require original sample records and assay metadata; accession pointers establish neither paired trajectories nor independent confirmation. Future animal-level prediction would require individual measurement\textendash{}outcome links, whereas inference from verified pools must remain at the pool level, without treating constituent cells or fields as independent animals.

The interaction-specific prior-art boundary also remains unassessed from primary and subsequent literature, so novelty is not established [C12, C13].

Confirmatory inference would require a separate pilot to estimate between-animal and between-clutch variation, justify a minimum-interest effect, and determine sample size before outcomes are examined. Response definitions, thresholds, allocation, uncertainty calculations, and stopping rules should be fixed before confirmation. Any predictor would then need evaluation on held-out biological units against stage-only, regeneration-initiating-cell/extracellular-matrix-only, and known-signaling baselines using matched information and prediction horizons, followed by independent biological confirmation.

No numerical biological effect with uncertainty, held-out improvement over those baselines, or independent functional result is established by the present evidence [C12, C13].

Explaining later recovery of competence requires an additional prediction and a verified later-competent X. laevis condition; it does not follow from a stage-40 versus stage-46 contrast, and a stage pattern from X. tropicalis cannot supply that missing observation. The evidence therefore changes the central argument from an output-timing hypothesis into a requirement for candidate-specific endogenous-state measurement: shortened memory becomes biologically testable only when persistence can be separated from continued input and measurement dynamics. This is a boundary on what would constitute a mechanistically discriminating explanation, not evidence that biological signal memory is absent [C11, C12, C13].

\section{Limitations}
\begin{itemize}
\item The original GEO sample records, expression files, and specimen mappings were not inspected. Applicability of the unlinked-observation result to the actual archival design is unknown.
\item The prior-art boundary rests on supplied summaries rather than inspected primary publications or a documented later-work search. No priority or novelty claim is supported.
\item The reporter construction, restricted crossed-pulse result, and accumulated-output symmetry are incompletely reproduced in the bounded preview. They require recovery of their full assumptions and derivations before being presented as completed proofs.
\item No named molecular intervening state or validated assay makes the shortened-memory candidate operational. The current conclusion is an evidence-readiness stop, not a biological refutation.
\item No empirical effect, confidence interval, predictive advantage, independent replication, or later regained-competence result is available.
\item Saved research session 44e759ac-38bb-4e8d-8c4c-4850541ad3dd ended with an incomplete review; call 18 (investigator) was not successfully completed. Earlier findings remain available, but this process limitation prevents claiming that the planned adversarial review was fully performed.
\end{itemize}
\section{Open obligations}
\begin{itemize}
\item Obtain original series and sample metadata for GSE245320, GSE245317, GSE245318, GSE245312, and GSE245313. Establish species, stage, time, tissue, animal versus pool, biological replication, library relationships, batch, assay units, and specimen-linked outcomes.
\item Determine whether early and late measurements are linked at the biological-unit level. If not, restrict archival inference to identifiable population contrasts unless additional identifying assumptions can be justified and tested.
\item Inspect the primary publications and subsequent literature through a specified investigation cutoff. Record species, exact interaction tested, experimental support, and source locations before assessing novelty.
\item Recover the complete reporter construction, including initial conditions and piecewise input intervals, and the complete restricted pulse-response and accumulated-output derivations. Check their assumptions and edge cases without treating reviewer agreement as independent verification.
\item Identify an endogenous basal-epidermal intervening state and supply candidate-specific localization, kinetic calibration, input-cessation, orthogonal-assay, and nonperturbation evidence before reopening the biological branch.
\item If archival analysis becomes feasible, define the biological replicate before modeling; address stage-batch confounding, pooling, homeolog mapping, module definitions, multiplicity, and tissue-composition changes.
\item Before predictive testing, preregister stage-only, regeneration-initiating-cell/extracellular-matrix-only, and known-signaling baselines; biological-unit validation splits; negative controls; the primary metric; and an uncertainty procedure.
\item For a future experiment, establish matched effective input profiles rather than relying only on delivered dose or unweighted exposure area. Include matched-time, sham, vehicle, wound-severity, viability, composition, and developmental-progression controls.
\item Use a separate pilot to justify a minimum useful effect or equivalence margin, estimate animal and clutch variation, and plan confirmatory precision. A confidence interval that still permits a meaningful effect is inconclusive.
\item Obtain independent unseen-specimen molecular and functional regeneration evidence, including a separately verified later-competent Xenopus laevis condition, before claiming an explanation of competence loss and recovery.
\end{itemize}
\section{Conclusions}
Unlinked stage-specific early and late distributions cannot, without additional identifying restrictions, establish an animal-level injury-to-epidermal coupling slope. This is a conditional mathematical limitation and does not establish the actual pairing structure of the GEO resource.

The shortened-memory hypothesis remains non-operational because its intervening molecular state and discriminating assay are unspecified. Conditional timing and reporter arguments help delimit the necessary observations but do not establish biological memory as the cause of regeneration failure.

The evidence supports neither a new mechanism nor a reproducible computational predictor. Progress requires inspectable primary records, candidate-specific measurement validation, and independent biological evidence rather than further interpretation of stage labels or endpoint persistence.

\section*{Claim and source notes}
These notes identify the saved evidence for claims used in the article. The complete claim register and reasoning are retained in the paper record.
\begin{sloppypar}
\small
\begin{description}
\item[C14] The substantive saved analysis consists of conditional identifiability constructions and simplified dynamical-model comparisons; it does not contain an executed transcriptomic reanalysis or biological experiment. \textit{Status: }established.
\textit{Evidence: }[Research session 21, artifact 2], [Research session 22, artifact 3], [Research session 22, artifact 5], [Research session 22, artifact 4]
\item[C1] The actual biological units, replicate counts, pooling, pairing, batches, assay units, and specimen-linked outcomes of GSE245320 and its named component series remain unverified in the supplied evidence. \textit{Status: }limitation.
\textit{Evidence: }[Research session 21, artifact 5], [Research session 22, artifact 4]
\item[C2] The project treats basal-epidermal regeneration-initiating cells and extracellular-matrix remodeling defects as known background, while the novelty of a stage-dependent timing or persistence interaction remains unassessed. \textit{Status: }limitation.
\textit{Evidence: }[Research session 21, artifact 1], [Research session 21, artifact 4], [Research session 22, artifact 2]
\item[C3] If observations contain only unlinked early and late stage-specific marginal distributions, without additional identifying restrictions, they do not identify the animal-level slope $\beta_s=\operatorname{Cov}(E,R\mid S=s)/\operatorname{Var}(E\mid S=s)$ when the denominator is positive. \textit{Status: }established.
\textit{Evidence: }[Research session 21, artifact 2], [Research session 21, artifact 5], [Research session 22, artifact 1]
\item[C4] A predictor defined solely by two developmental-stage means cannot supply an independently estimable interaction with the binary stage indicator in a model already containing stage. \textit{Status: }established.
\textit{Evidence: }[Research session 22, artifact 5]
\item[C5] The saved first-order reporter proof candidate states that calibrated readings at $0.5$ and $3$ hours can be identical for sustained and interrupted nonnegative inputs, even when reporter gain and decay are known. \textit{Status: }unresolved.
\textit{Evidence: }[Research session 22, artifact 3], [Research session 22, artifact 4]
\item[C6] A candidate explanation is shorter basal-epidermal injury-signal memory at stage 46 than at stage 40, predicting a larger relative transcriptional benefit from shifting a matched input pulse nearer the regenerative response window. \textit{Status: }conjecture.
\textit{Evidence: }[Research session 22, artifact 2], [Research session 22, artifact 6]
\item[C7] At a single terminal readout, an exponential signal-memory response is observationally equivalent to a common-memory response with pulse-time-dependent inducibility. \textit{Status: }established.
\textit{Evidence: }[Research session 22, artifact 2]
\item[C8] The saved restricted-model analysis proposes that crossing two pulse times with two post-pulse delays can distinguish shortened memory from pulse-time-only inducibility, while an overlapping absolute readout time additionally tests readout-time-only wound-clock gating. \textit{Status: }unresolved.
\textit{Evidence: }[Research session 22, artifact 5]
\item[C9] Allowing both pulse-time-dependent encoding and readout-time-dependent gating can reproduce an exponential memory response across pulse times and delays. \textit{Status: }established.
\textit{Evidence: }[Research session 22, artifact 5]
\item[C10] The saved accumulated-output model argument reports that exchanging signal-state decay and output turnover leaves its impulse response unchanged, so accumulated transcription alone cannot generally identify which process is faster. \textit{Status: }unresolved.
\textit{Evidence: }[Research session 22, artifact 5], [Research session 22, artifact 6]
\item[C11] The shortened-memory hypothesis is presently non-operational because no endogenous basal-epidermal intervening state or candidate-specific, stage-calibrated assay is supplied that separates its decay from continued input and measurement filtering. \textit{Status: }limitation.
\textit{Evidence: }[Research session 22, artifact 5], [Research session 22, artifact 6]
\item[C12] No numerical biological effect with uncertainty, improvement over known-mechanism baselines, independent biological confirmation, or evidence explaining later recovery of competence is established. \textit{Status: }limitation.
\textit{Evidence: }[Research session 22, artifact 4], [Research session 22, artifact 6], [Research session 21, artifact 1]
\item[C13] A future discriminating test would require a named intervening state, verified input cessation, stage-specific measurement calibration, and an orthogonal endogenous-state assay; a precise result excluding the prespecified faster-decay effect would oppose the designated hypothesis, whereas an imprecise null would be inconclusive. \textit{Status: }unresolved.
\textit{Evidence: }[Research session 22, artifact 6], [Research session 22, artifact 4], [Research session 21, artifact 1], [Research session 21, artifact 5]
\end{description}
\end{sloppypar}
\section*{Paper writing source material used}
These sources informed manuscript organization and writing methods, not the scientific evidence or findings.
\begin{itemize}
\item PLOS ONE \textemdash{} Manuscript and Reporting Guidelines (general): \url{https://journals.plos.org/plosone/s/submission-guidelines}
\item EASE \textemdash{} Guidelines for Authors of Scientific Articles (full text) (general): \url{https://pmc.ncbi.nlm.nih.gov/articles/PMC4130677/}
\end{itemize}
\section*{Authorship note}
AI-assisted synthesis from saved research; not independently peer reviewed.
\end{document}
